% practice.tex -- Module 7 (The Nature of Light): Practice set.
% 23 new HSC-style questions (8 multiple choice, 10 short answer, 5 extended response),
% ordered from easy to HSC level, covering Chapters 1-4.
%   (solve.py, bbcheck.py, verify_answers.py; outputs saved next to them).
% No homework, booklet or exam question (or its numbers) is reused.
% Answers: practice_answers.tex (one practiceans per practiceq, same order).
% \Needspace before each question keeps a question (and its options or parts) on one
% page when it fits; it is conditional and never forces a page break by itself.
\unnumberedchapter[\faDumbbell]{Practice set}

% Marks for one part of a question, pushed to the right margin; if they do not fit on
% the last line, they move to a new line as a whole (TeXbook \signed trick).
% New macro for the practice set only (tell the design-system owner).
\providecommand\practiceMarks[1]{{\unskip\nobreak\hfil\penalty50\hskip1em\hbox{}\nobreak
  \hfill\mbox{\small\textcolor{hscmuted}{(#1)}}\parfillskip=0pt\finalhyphendemerits=0\par}}

% Small black body sketches for Question 4 (Planck curves; x = wavelength in um,
% normalised so that the 6000 K curve peaks at 1).
\pgfplotsset{pqbbsketch/.style={hsc, width=64mm, height=44mm,
    xmin=0, xmax=2.6, ymin=0, ymax=1.3, xtick=\empty, ytick=\empty, grid=none,
    axis lines=left, xlabel={wavelength}, ylabel={intensity},
    label style={font=\footnotesize},
    xlabel style={at={(axis description cs:0.5,0)}, anchor=north, yshift=-1pt},
    ylabel style={at={(axis description cs:0,0.5)}, anchor=south, yshift=1pt},
    legend style={at={(1,1)}, anchor=north east, font=\scriptsize, draw=none,
      fill=white, inner sep=1pt, row sep=-1pt},
    legend cell align=left,
    domain=0.10:2.55, samples=160, smooth,
    declare function={pqbb(\x,\t)=(\x)^(-5)/(exp(14388/(\x*\t))-1)/0.267374;}}}

These 23 questions cover all four chapters. They start easy and finish at HSC level. The
\textbf{Ch} tag next to each question number shows the chapter it tests. Click the tag to go
to that chapter.
\begin{itemize}
  \item Total: 71 marks. Allow about 2 hours. (In the HSC exam you have about 1.8 minutes for
        each mark.)
  \item Use only a calculator and the HSC data sheet and formulae sheet.
  \item Show all your working. Give every answer with a unit and sensible significant figures.
  \item Full worked answers and marking criteria are in \emph{Practice set answers}, at the end
        of the Answers chapter. Click a question number to jump to its answer.
\end{itemize}

%% ==================================================================
\section*{Section I: multiple choice (8 marks)}
Choose the best answer: A, B, C or D.

\Needspace{11\baselineskip}
\begin{practiceq}{1}{1}\label{pq:produce}
Which of these produces an electromagnetic wave?
\choices[1]{a charged metal sphere at rest on an insulating stand}
  {an electron moving at a constant velocity in a vacuum}
  {electrons moving up and down, again and again, in a radio antenna}
  {a steady direct current in a long straight wire}
\end{practiceq}

\Needspace{7\baselineskip}
\begin{practiceq}{1}{4}\label{pq:laser}
A spacecraft moves towards a space station at $0.80c$. It sends a laser pulse forward, towards
the station. What speed does an observer on the station measure for the laser pulse?
\choices{$0.20c$}{$1.00c$}{$1.67c$}{$1.80c$}
\end{practiceq}

\Needspace{10\baselineskip}
\begin{practiceq}{1}{2}\label{pq:foucault}
In 1850 Foucault measured the speed of light in water. Light travelled more slowly in water
than in air. What did this result show?
\choices[1]{Newton's particle explanation of refraction was wrong.}
  {Light is an electromagnetic wave.}
  {Light is a transverse wave.}
  {Light is a stream of photons.}
\end{practiceq}

\Needspace{25\baselineskip}
\begin{practiceq}{1}{3}\label{pq:bbgraph}
Which graph best shows the black body radiation curves of a star at \qty{4000}{K} and a star
at \qty{6000}{K}?
\choices[2]{%
  \begin{tikzpicture}[baseline=(current bounding box.north)]
    \begin{axis}[pqbbsketch]
      \addplot[hscseriesA, line width=1.1pt, domain=0:2.55] {exp(-((x-0.483)/0.22)^2)};
      \addplot[hscseriesB, line width=1.1pt, dashed, domain=0:2.55] {0.132*exp(-((x-0.7245)/0.33)^2)};
      \legend{\qty{6000}{K}, \qty{4000}{K}}
    \end{axis}
  \end{tikzpicture}}
  {%
  \begin{tikzpicture}[baseline=(current bounding box.north)]
    \begin{axis}[pqbbsketch]
      \addplot[hscseriesA, line width=1.1pt] {7.59375*pqbb(x,4000)};
      \addplot[hscseriesB, line width=1.1pt, dashed] {pqbb(x,6000)/7.59375};
      \legend{\qty{6000}{K}, \qty{4000}{K}}
    \end{axis}
  \end{tikzpicture}}
  {%
  \begin{tikzpicture}[baseline=(current bounding box.north)]
    \begin{axis}[pqbbsketch]
      \addplot[hscseriesA, line width=1.1pt] {pqbb(x,6000)};
      \addplot[hscseriesB, line width=1.1pt, dashed] {4*pqbb(x,4000)};
      \legend{\qty{6000}{K}, \qty{4000}{K}}
    \end{axis}
  \end{tikzpicture}}
  {%
  \begin{tikzpicture}[baseline=(current bounding box.north)]
    \begin{axis}[pqbbsketch]
      \addplot[hscseriesA, line width=1.1pt] {pqbb(x,6000)};
      \addplot[hscseriesB, line width=1.1pt, dashed] {pqbb(x,4000)};
      \legend{\qty{6000}{K}, \qty{4000}{K}}
    \end{axis}
  \end{tikzpicture}}
\end{practiceq}

\Needspace{13\baselineskip}
\begin{practiceq}{1}{1}\label{pq:starlines}
Every hydrogen line in the spectrum of star P is at a slightly longer wavelength than the same
line in a laboratory spectrum. The lines are also much wider than the same lines in the Sun's
spectrum. Which conclusion do these observations support best?
\choices[1]{P is moving towards Earth, and its atmosphere has a lower density than the Sun's.}
  {P is moving away from Earth, and its surface is hotter than the Sun's.}
  {P is moving towards Earth, and it rotates faster than the Sun.}
  {P is moving away from Earth, and it rotates faster or has a denser atmosphere than the Sun.}
\end{practiceq}

\Needspace{7\baselineskip}
\begin{practiceq}{1}{2}\label{pq:twofilters}
Unpolarised light of intensity \qty{360}{W.m^{-2}} passes through two polarising filters. The
angle between their transmission axes is \ang{50.0}. What is the intensity of the light that
leaves the second filter?
\choices{\qty{74.4}{W.m^{-2}}}{\qty{106}{W.m^{-2}}}{\qty{116}{W.m^{-2}}}{\qty{149}{W.m^{-2}}}
\end{practiceq}

\Needspace{11\baselineskip}
\begin{practiceq}{1}{3}\label{pq:lasers}
Laser R gives out red light (\qty{660}{nm}). Laser B gives out blue light (\qty{470}{nm}). Both
lasers have the same power. Which statement is correct?
\choices[1]{Each photon from laser B has less energy than each photon from laser R.}
  {Photons from laser B travel faster than photons from laser R.}
  {Laser R gives out more photons each second than laser B.}
  {The beam from laser B carries more energy each second than the beam from laser R.}
\end{practiceq}

\Needspace{11\baselineskip}
\begin{practiceq}{1}{4}\label{pq:proper}
A spacecraft travels at a constant velocity from Earth to a space station. The station is at
rest relative to Earth. Who measures the proper time and the proper length for this trip?
\choices[1]{Proper time: clocks on Earth. \quad Proper length: the crew.}
  {Proper time: the crew's clock. \quad Proper length: observers on Earth.}
  {Both the proper time and the proper length: the crew.}
  {Both the proper time and the proper length: observers on Earth.}
\end{practiceq}

%% ==================================================================
\section*{Section II: short answer (33 marks)}
Show all working for calculations.

\Needspace{12\baselineskip}
\begin{practiceq}{3}{1}\label{pq:microwave}
A student takes the turntable out of a microwave oven. The oven's label gives a frequency of
\qty{2.45}{GHz}. She heats a flat bar of chocolate until a line of small spots melts.
Neighbouring melted spots are half a wavelength apart. The centres of the first and the fifth
melted spots are \qty{25.0}{cm} apart.
\begin{enumerate}[label=(\alph*)]
  \item Calculate the speed of light from these results. \practiceMarks{2 marks}
  \item Suggest one change to the method that would improve the accuracy or the reliability of
        the result. Give a reason. \practiceMarks{1 mark}
\end{enumerate}
\end{practiceq}

\Needspace{9\baselineskip}
\begin{practiceq}{3}{3}\label{pq:filament}
The tungsten filament of a lamp is at \qty{2700}{\degreeCelsius}. Treat it as a black body.
\begin{enumerate}[label=(\alph*)]
  \item Calculate the peak wavelength of its radiation. Name the part of the electromagnetic
        spectrum it is in. \practiceMarks{2 marks}
  \item Calculate the energy, in electron volts, of one photon at this peak wavelength.
        \practiceMarks{1 mark}
\end{enumerate}
\end{practiceq}

\Needspace{9\baselineskip}
\begin{practiceq}{3}{4}\label{pq:pions}
Charged pions are unstable particles. At rest, their mean lifetime is \qty{2.60e-8}{s}. In a
beam from an accelerator, pions travel at $0.920c$ relative to the laboratory.
\begin{enumerate}[label=(\alph*)]
  \item Calculate the mean lifetime of these pions, measured in the laboratory.
        \practiceMarks{2 marks}
  \item Calculate the average distance a pion travels in the laboratory before it decays.
        \practiceMarks{1 mark}
\end{enumerate}
\end{practiceq}

\Needspace{11\baselineskip}
\begin{practiceq}{3}{2}\label{pq:thirty}
Unpolarised light passes through polariser P1, which has a vertical transmission axis. It then
passes through polariser P2. The light that leaves P2 has 30\,\% of the intensity of the
unpolarised light.
\begin{enumerate}[label=(\alph*)]
  \item Calculate the angle between the transmission axes of P1 and P2. \practiceMarks{2 marks}
  \item A third polariser, P3, with a horizontal transmission axis, is placed after P2. What
        percentage of the original intensity leaves P3? \practiceMarks{1 mark}
\end{enumerate}
\end{practiceq}

\Needspace{7\baselineskip}
\begin{practiceq}{3}{2}\label{pq:doubleslit}
Laser light passes through two narrow slits \qty{0.350}{mm} apart. Bright fringes form on a
screen \qty{2.60}{m} away. The first-order ($m = 1$) and the sixth-order ($m = 6$) bright
fringes are on the same side of the central fringe. Their centres are \qty{23.5}{mm} apart.
Calculate the wavelength of the laser light.
\end{practiceq}

\Needspace{10\baselineskip}
\begin{practiceq}{4}{3}\label{pq:magnesium}
Ultraviolet light of wavelength \qty{220}{nm} falls on a clean magnesium surface. The work
function of magnesium is \qty{3.66}{eV}. Calculate:
\begin{enumerate}[label=(\alph*)]
  \item the energy of each photon, in eV \practiceMarks{1 mark}
  \item the maximum kinetic energy of the photoelectrons, in eV \practiceMarks{1 mark}
  \item the stopping voltage \practiceMarks{1 mark}
  \item the maximum speed of the photoelectrons. \practiceMarks{1 mark}
\end{enumerate}
\end{practiceq}

\Needspace{9\baselineskip}
\begin{practiceq}{3}{2}\label{pq:grating}
Light from a hydrogen discharge tube passes through a diffraction grating with 420 lines per
mm. The spectrum contains a red line of wavelength \qty{656}{nm}.
\begin{enumerate}[label=(\alph*)]
  \item Calculate the angle of the first-order red line. \practiceMarks{2 marks}
  \item What is the highest order in which this red line can be seen? Justify your answer.
        \practiceMarks{1 mark}
\end{enumerate}
\end{practiceq}

\Needspace{10\baselineskip}
\begin{practiceq}{3}{4}\label{pq:muonpair}
A muon is a particle with a rest mass of \qty{1.884e-28}{kg}. Its antiparticle is the antimuon,
which has the same mass. A muon and an antimuon, both almost at rest, annihilate. They make two
gamma-ray photons.
\begin{enumerate}[label=(\alph*)]
  \item Calculate the total energy released. \practiceMarks{1 mark}
  \item Calculate the energy of each photon, in MeV. \practiceMarks{1 mark}
  \item Explain why two photons are made, not one. \practiceMarks{1 mark}
\end{enumerate}
\end{practiceq}

\Needspace{12\baselineskip}
\begin{practiceq}{4}{3}\label{pq:planck}
Classical physics could explain the long-wavelength part of a black body curve, but not the
short-wavelength part.
\begin{enumerate}[label=(\alph*)]
  \item Describe what classical physics predicted for the intensity at short wavelengths.
        \practiceMarks{1 mark}
  \item Use Planck's hypothesis to explain why the measured intensity falls to zero at short
        wavelengths. \practiceMarks{2 marks}
  \item State one way in which Einstein's photon model went further than Planck's hypothesis.
        \practiceMarks{1 mark}
\end{enumerate}
\end{practiceq}

\Needspace{9\baselineskip}
\begin{practiceq}{4}{4}\label{pq:momentum}
An electron in a medical linear accelerator moves at \qty{2.79e8}{m.s^{-1}}.
\begin{enumerate}[label=(\alph*)]
  \item Calculate its relativistic momentum. Compare your answer with $m_0v$. \practiceMarks{2 marks}
  \item The accelerator is made longer, so it does more work on the electron. Explain why the
        electron's speed can never reach \qty{3.00e8}{m.s^{-1}}. \practiceMarks{2 marks}
\end{enumerate}
\end{practiceq}

%% ==================================================================
\section*{Section III: extended response (30 marks)}
Plan your answer before you write. Use correct physics words, and link each idea to the next.

\Needspace{5\baselineskip}
\begin{practiceq}{5}{1}\label{pq:maxwell}
Analyse how Maxwell's theory of electromagnetism changed scientists' understanding of light.
Include a calculation that uses values from the data sheet.
\end{practiceq}

\Needspace{24\baselineskip}
\begin{practiceq}{6}{1}\label{pq:stars}
An astronomer records the spectra of two stars, A and B. The table compares their absorption
lines with laboratory lines. It also gives the peak wavelength of each star's continuous
spectrum.
\begin{center}\small
  \begin{tabular}{@{}L{54mm}C{27mm}C{27mm}C{27mm}@{}}
    \toprule
    \colhead{Absorption line} & \colhead{Laboratory (\si{nm})} & \colhead{Star A (\si{nm})} &
      \colhead{Star B (\si{nm})}\\
    \midrule
    Hydrogen & 434.05 & 433.92 & 434.15\\
    Hydrogen & 486.13 & 485.98 & 486.24\\
    Sodium & 589.00 & 588.82 & not seen\\
    Sodium & 589.59 & 589.41 & not seen\\
    Hydrogen & 656.28 & 656.08 & 656.43\\
    \midrule
    Width of the lines & & narrow & very broad\\
    Peak of the continuous spectrum & & \qty{725}{nm} & \qty{322}{nm}\\
    \bottomrule
  \end{tabular}
\end{center}
\begin{enumerate}[label=(\alph*)]
  \item Calculate the surface temperature of each star. \practiceMarks{2 marks}
  \item Analyse what the absorption lines show about the chemical composition, the motion,
        and the rotation or density of each star. \practiceMarks{4 marks}
\end{enumerate}
\end{practiceq}

\Needspace{6\baselineskip}
\begin{practiceq}{6}{2}\label{pq:models}
In the 1700s most scientists used Newton's particle model of light. By the late 1800s they used
a model of light as a transverse wave. Analyse how experimental evidence about refraction,
interference and polarisation led to this change.
\end{practiceq}

\Needspace{13\baselineskip}
\begin{practiceq}{6}{4}\label{pq:muons}
Cosmic rays make muons \qty{9.00}{km} above sea level (measured in the Earth frame). The muons
travel straight down at $0.9960c$. The mean lifetime of a muon at rest is
\qty{2.20}{\micro\second}.
\begin{enumerate}[label=(\alph*)]
  \item Calculate how far a muon would travel in one mean lifetime if there were no time
        dilation. \practiceMarks{1 mark}
  \item In the Earth frame, calculate the mean lifetime of these muons and the time they take
        to reach sea level. \practiceMarks{2 marks}
  \item In the muon's frame, calculate the thickness of the \qty{9.00}{km} of atmosphere.
        \practiceMarks{1 mark}
  \item Detectors at sea level find that a large fraction of these muons arrive. Explain how
        this supports special relativity. Refer to your answers to (a)--(c). \practiceMarks{2 marks}
\end{enumerate}
\end{practiceq}

\Needspace{48\baselineskip}
\begin{practiceq}{7}{3}\label{pq:pegraph}
A student shines ultraviolet light of different wavelengths onto a metal plate inside a vacuum
photocell. For each wavelength she measures the stopping voltage $V_0$.
\begin{center}\small
  \begin{tabular}{@{}L{44mm}*{6}{C{13mm}}@{}}
    \toprule
    Wavelength $\lambda$ (\si{nm}) & 240 & 265 & 290 & 315 & 340 & 365\\
    Frequency $f$ ($\times10^{14}$ \si{Hz})\rule{0pt}{7mm} & \rule{9mm}{0.4pt} & \rule{9mm}{0.4pt} &
      \rule{9mm}{0.4pt} & \rule{9mm}{0.4pt} & \rule{9mm}{0.4pt} & \rule{9mm}{0.4pt}\\
    Stopping voltage $V_0$ (\si{V}) & 2.33 & 1.77 & 1.43 & 1.05 & 0.81 & 0.52\\
    \bottomrule
  \end{tabular}
\end{center}
\begin{enumerate}[label=(\alph*)]
  \item Complete the table by calculating the frequency for each wavelength. \practiceMarks{1 mark}
  \item On the grid, plot a graph of $V_0$ (vertical axis) against $f$ (horizontal axis).
        Draw a line of best fit. \practiceMarks{2 marks}
  \item Use the gradient of your line to find a value for Planck's constant. \practiceMarks{2 marks}
  \item Use your graph to find the threshold frequency and the work function of the metal.
        (Extend your line if you need to.) Then use the table below to identify the metal.
        \practiceMarks{2 marks}
\end{enumerate}
\begin{center}\small
  \begin{tabular}{@{}L{36mm}C{20mm}C{20mm}C{20mm}C{20mm}@{}}
    \toprule
    Metal & potassium & calcium & magnesium & zinc\\
    Work function (\si{eV}) & 2.29 & 2.87 & 3.66 & 4.30\\
    \bottomrule
  \end{tabular}
\end{center}
\begin{center}
  \begin{tikzpicture}[x=1mm, y=1mm]
    \draw[step=2, hscrule!70, line width=0.2pt] (0,0) grid (130,80);
    \draw[step=10, hscmuted!45, line width=0.45pt] (0,0) grid (130,80);
  \end{tikzpicture}
\end{center}
\end{practiceq}
